\documentclass[11pt]{article}
\usepackage[a4paper,margin=1in]{geometry}
\usepackage{amsmath,amssymb,amsthm}
\usepackage{booktabs}
\usepackage[hidelinks]{hyperref}
\usepackage{natbib}

\title{Numerical Optimization of Penalized Trend Smoothness\\
by Pooled Leave-One-Out Cross-Validation}
\author{Heriberto Espino Montelongo\\[3pt]\small Working draft --- author affiliations and coauthors to confirm}
\date{18 October 2026 --- research outline, not a submission draft}

\begin{document}
\maketitle

\begin{abstract}
Penalized least-squares trend smoothing requires a choice of its
regularization parameter. Following the experimental setting of
Cort\'es-Toto, Guerrero and Reyes (2017), we investigate the numerical
optimization of leave-one-out cross-validation (LOOCV) loss for one
or more comparable observed time series. We define a pooled criterion,
derive its finite-parameter derivative and exact endpoint limits,
and specify comparisons between bounded scalar minimization,
multi-bracket Brent root refinement and restricted exact Sturm
root isolation. A completed source-inspired pilot is complemented by a separate
stochastic-trend simulation comparing observed-series LOOCV against
the known noise-variance ratio in Guerrero's model. A financial
illustration remains a prospective application. No accuracy or
efficiency superiority result is asserted at this stage.
\end{abstract}

\section{Introduction and literature positioning}
Trend smoothing by penalized least squares (PLS) is classical
\citep{guerrero2007,cortes2017}. Classical LOOCV and GCV
are already used for choosing smoothing parameters
\citep{cortes2017,weinert2007}.
A new numerical contribution requires a demonstrated improvement
over existing global GCV optimization, including the prior-art
investigation discussed in the research notebook.
We address observed-series reconstruction, not prediction.

\section{Pooled LOOCV and normalized smoothness}
Let $y^{(r)}\in\mathbb R^N$ for $r=1,\dots,R$, with
$N>d\geq1$, $D_d$ the forward difference matrix,
$Q=D_d^\top D_d$, and
\begin{equation}
 H_\lambda=(I_N+\lambda Q)^{-1}, \qquad
 \widehat\tau_\lambda^{(r)}=H_\lambda y^{(r)}.
\end{equation}
For $\lambda>0$, define the pooled LOOCV MSE
\begin{equation}\label{eq:pooled}
 F_R(\lambda)=\frac1{RN}\sum_{r=1}^{R}\sum_{i=1}^{N}
 \left(\frac{y_i^{(r)}-(H_\lambda y^{(r)})_i}
 {1-(H_\lambda)_{ii}}\right)^2.
\end{equation}
This averages errors \emph{before} minimizing;
it is not the average of $R$ optimum regularization parameters.
The normalized Guerrero-type coordinate is
\begin{equation}
 S(\lambda)=\frac{N-\operatorname{tr}H_\lambda}{N-d}
 \in[0,1].
\end{equation}
For $R=1$, Eq.~\eqref{eq:pooled} is ordinary LOOCV.

\section{Endpoint limits and analytical derivatives}
Put $T_\lambda=Q(I+\lambda Q)^{-1}$. Since
$I-H_\lambda=\lambda T_\lambda$, LOOCV residuals satisfy
\begin{equation}
 e_{ri}(\lambda)=\frac{(T_\lambda y^{(r)})_i}
 {(T_\lambda)_{ii}} .
\end{equation}
The continuous extension at $\lambda=0$ is
$e_{ri}(0)=(Qy^{(r)})_i/Q_{ii}$. At
$\lambda=\infty$, write $P$ for the orthogonal projector onto
$\ker Q$ and use
$e_{ri}(\infty)=[(I-P)y^{(r)}]_i/(1-P_{ii})$.
For finite $\lambda$, $T'_\lambda=-T_\lambda^2$ and hence
\begin{equation}
 F_R'(\lambda)=\frac{2}{RN}\sum_{r,i}e_{ri}(\lambda)
 \frac{-(T_\lambda^2y^{(r)})_i(T_\lambda)_{ii}
 +(T_\lambda y^{(r)})_i(T_\lambda^2)_{ii}}
 {(T_\lambda)_{ii}^2}.
\end{equation}
These are standard linear-smoother calculus identities; proofs
and numerical edge cases are recorded in notes/PROOF_LEDGER.md.

\section{Numerical candidate recovery}
We plan to compare (i) a finite dense log-$\lambda$ grid,
(ii) one bounded Brent-style scalar search,
(iii) multiple sign-change brackets of the derivative followed
by Brent root refinement and evaluation of both endpoints,
and (iv) exact Sturm isolation for small rational instances.
Only method (iv), subject to exact-arithmetic restrictions and
correct root classification, supplies exact positive-root counts.
A derivative-sign scan alone cannot guarantee completeness.

\section{Simulation design and financial illustration}
Use the four-factor design of \citet{cortes2017}:
linear or nonlinear Beta-mixture trend, seasonality present/absent,
Gaussian noise standard deviation 0.5/2, and $N=50/200$.
Pools with $R>1$ are an independently declared extension
using a common latent trend and independent noisy replicates.
Predeclare paired comparisons of candidate recovery, objective
values, runtime, boundary selection and sensitivity to pooling.
A financial case study may illustrate retrospective smoothing,
without forecasting or investment-performance claims.

\subsection{Known noise-variance ratio under Guerrero's model}

To distinguish numerical LOOCV optimization from model-informed trend
estimation, we introduce a \emph{separate} stochastic design following
\citet{guerrero2007}. For mutually uncorrelated disturbances, let
\begin{equation}
 y^{(r)}=\tau^{(r)}+\eta^{(r)},\quad
 D_d\tau^{(r)}=\varepsilon^{(r)},\quad
 \operatorname{Var}(\eta^{(r)})=\sigma_\eta^2I,\quad
 \operatorname{Var}(\varepsilon^{(r)})=\sigma_\varepsilon^2I.
\end{equation}
The population variance ratio motivates the model-informed parameter
\begin{equation}
 \lambda_G=\frac{\sigma_\eta^2}{\sigma_\varepsilon^2}.
\end{equation}
The benchmark is given \emph{true simulation variances}, not estimated
variance components from empirical financial observations. We compare
$\lambda_G$ with the pooled LOOCV selections produced by one bounded
Brent run, multibracket derivative Brent, and a finite dense grid.

For each choice $\lambda_m$ we measure (i) observed-series
$F_R(\lambda_m)$ and (ii) the latent-trend reconstruction criterion
\begin{equation}
 \mathrm{TrendMSE}(\lambda_m)=
 \frac{1}{RN}\sum_{r=1}^R
 \|H_{\lambda_m}y^{(r)}-\tau^{(r)}\|_2^2.
\end{equation}
The simulated latent trend is \emph{never} used to select $\lambda$
under LOOCV. The optimization objective and the reconstruction target
are different; obtaining the smallest observed-sample CV criterion
does not establish the smallest realized or expected latent-trend
error. The fixed polynomial initialization of the integrated trend
should not be conflated with a proper finite-variance prior for
its unpenalized nullspace coefficients.

\subsection{Exploratory known-ratio pilot (not confirmatory)}
A separate stochastic-trend pilot completed $12$ settings,
each with five randomly seeded repetitions and pool sizes $R=1$
and $R=3$, yielding $120$ nested pool evaluations and four method
rows per evaluation. The two pool sizes share the first generated
series at matching seeds; hence they cannot be treated as
independent samples in a pooled inferential analysis.

For the $R=3$ group (60 pool evaluations), the mean observed
LOOCV scores were $2.697581$ for the known Guerrero variance
ratio and $2.681409$ for the multi-bracket LOOCV selection.
The mean \emph{true trend} MSEs were instead $0.385673$
and $0.394848$, respectively. In $45$ of those $60$
paired pools the Guerrero ratio attained lower realized
trend MSE. This is a known-parameter comparison;
the advantage cannot be attributed to an estimated variance
component method.

In this stochastic pilot, both bounded Brent and
multi-bracket Brent attained effectively identical LOOCV
objectives in every pool at a tolerance of $10^{-7}$.
All $60$ $R=3$ objectives yielded only one
numerically detected derivative-sign root under the
finite scan. This count is \emph{not} a global certificate
or a proof of unimodality. Therefore, this particular
DGP provides evidence for the distinction between
statistical and numerical targets, but no measured
numerical superiority of multi-bracket search over the
single-start procedure. The results require a genuinely
independent confirmatory experiment and objective-specific
numerical verification before journal-level claims.

\subsection{Feasible estimation of both variance components: exploratory pilot}

The preceding known-ratio comparison assumes the true variance
components are supplied. This is not a feasible smoothing selector
for ordinary data. Under the additional independent Gaussian
stochastic-difference model and $V=I$, put $z^{(r)}=D_2y^{(r)}$ and
$A=D_2D_2^\top$. Then
\[
 \operatorname{Cov}(z^{(r)})=
 \sigma_\varepsilon^2 I+\sigma_\eta^2 A.
\]
We independently implement two \emph{feasible extensions}: (a)
Gaussian maximum likelihood of the differenced observations, with
the common covariance scale profiled analytically, and (b) two
quadratic moment equations with nonnegative least-squares
constraints. These estimators are \emph{not} attributed to
Guerrero's published smoothness-selection procedure.

A distinct exploratory pilot considered $N\in\{50,200\}$,
$\sigma_\eta\in\{0.5,2\}$,
$\sigma_\varepsilon\in\{0.1,0.3,1\}$,
five independently seeded pools per cell, and nested $R=1$ and $R=3$.
For the primary $R=3$ subset (60 pools), the mean
observed-series LOOCV and latent-trend MSE results were:
\begin{center}
\begin{tabular}{lrr}
\toprule
Parameter selection & LOOCV MSE & Latent trend MSE\\
\midrule
Known true variance ratio (oracle) & 2.697581 & \textbf{0.385673}\\
Estimated Gaussian difference likelihood & 2.696789 & 0.387916\\
Pooled LOOCV (bounded Brent) & \textbf{2.681409} & 0.394848\\
Estimated nonnegative quadratic moments & 102.033035 & 96.127631\\
\bottomrule
\end{tabular}
\end{center}
Gaussian likelihood achieved lower realized trend MSE than bounded
LOOCV in 37 of 60 $R=3$ paired pools. It estimated observation-noise
and trend-innovation variance with average relative absolute errors
of approximately 10.0\% and 18.3\%, respectively. The error rates
fell with increasing sample length in this particular simulation.

The constrained moment method set its estimated trend-innovation
variance to zero in 16 of 60 $R=3$ pools, implying
$\widehat\lambda=\infty$ and fitting a linear polynomial even
when the stochastic trend had nonzero second differences. Such
catastrophic errors make its mean latent-trend MSE unreliable as a
measure of typical numerical approximation accuracy; we report
failure rates and medians as well. The comparison is necessarily
exploratory with just five seeds per DGP cell, and results
do not establish an expected-risk inequality for all noise
configurations. The $R=1$ observations are nested in $R=3$,
not independent replications. Complete per-cell results and
replication details are retained in the corresponding notes.

\section{Discussion and limitations}
\textbf{Confirmatory results pending.} The exploratory pilots have been recorded, but neither numerical dominance nor a
novel global-solver theorem is established. Related algorithms,
correlated errors, scalar minimization budgets and
computational exactness require further investigation.

\bibliographystyle{plainnat}
\bibliography{references}
\end{document}
